% =====================================================================
%  Цель и задание лабораторной работы.
% =====================================================================

\textbf{Цель:} разработать методику комплексной оценки защищённости
сторонних программных систем и применить её на практике к реальному
приложению, действуя со стороны внешнего наблюдателя («чёрный ящик»), а
не со стороны разработчика системы.

\vspace{1em}

\textbf{Задание}:
\begin{enumerate}
  \item Разработать методику оценки защищённости систем.
  \item Выполнить анализ предоставленного приложения с позиций: политики
  безопасности операционной системы; политики безопасности компонентов
  приложения и их взаимодействия между собой; сетевой политики
  безопасности; безопасности конфиденциальности данных.
  \item Предложить комплекс мер по усовершенствованию приложения.
  \item Предложить рекомендации по безопасному использованию приложения.
  \item Оформить отчёт о проделанной работе.
\end{enumerate}

\textbf{Постановка задачи.} В качестве «предоставленного приложения»
(инструментальное средство №2 задания) выступает вариант, выполненный в
рамках лабораторной работы №2.1, -- веб-приложение \texttt{SecureApp}
(\texttt{Lab2.1-SecureApp}). Разница с лабораторной работой №2.1
принципиальна: там анализ велся от лица разработчика, знающего исходный
код и намеренно закладывающего в него конкретные защитные механизмы;
здесь тот же набор направлений анализа (ОС, компоненты и их
взаимодействие, сеть, конфиденциальность) выполняется от лица внешнего
аудитора -- с помощью чёрного ящика и целенаправленных проверок, --
а исходный код привлекается только для подтверждения или опровержения
гипотез, выдвинутых по результатам чёрного ящика, а не как основной
источник информации.
